%%%%%%%%%%%%%%%%%%%%%%%%%%%%%%%%%%%%%%%%%%%%%%%%%%%%%%%%%%%%%%%%%%%%%
%%% 
%%% Einstellungen
%%%
%%% Benoetigt wird KOMA-Skript ab Version 2.8j vom 30.07.2001
%%% Vor der Veraenderung irgendwelcher Einstellungen wird dringend empfohlen
%%% die Anleitung zum KOMA-Skript (scrguide) zu konsultieren !!
%%%
%%% Bilder bitte nicht mit Endung einbinden, so ist die Erzeugung von
%%% DVI, PS und PDF problemlos möglich!
%%% Infos zur \documentclass:
%%% http://wannert.net/latex/eine-praeambel-nicht-nur-fuer-mathematiker
%%%
%%%%%%%%%%%%%%%%%%%%%%%%%%%%%%%%%%%%%%%%%%%%%%%%%%%%%%%%%%%%%%%%%%%%%

% Haupteinstellungen für die KOMA-Dokumentenvorlage
% ===================================================
\documentclass[    
            paper=a4, 		%Papiergröße
            %----------------------------------------------
            %DIVcalc,        % Automatische Berechnung des Satzspiegels nach "goldenem Schnitt"
            DIV15,         % Unterteilung des Seitenlayouts in xy Streifen,
                           % führt zu kleineren Rändern, je größer die Zahl.
            %----------------------------------------------
            oneside,        % Seitenaufteilung: einteilig.
            %twoside,         % Seitenaufteilung: zweiseitig.
            %----------------------------------------------
            titlepage,      % Titelseite speziell, ohne Seitenkopf und -fu{\ss}
            %----------------------------------------------
            fontsize=12pt,           % Schriftgröße
            %----------------------------------------------
            %openright,      % Kapitel beginnen nur auf rechten Seiten.
            openany, 				% Kapitel können auf geraden und ungeraden Seiten beginnen
            %----------------------------------------------
            headinclude,    % Kopfzeile wird Seiten-Layouts mit beruecksichtigt
            %----------------------------------------------
            headsepline,    % Linie unter Kolumnentiteln an
            %headnosepline,  % Linie unter Kolumnentiteln aus
            %----------------------------------------------
            %footexclude,    % Fu{\ss}noten-Bereich ausschalten?
            footinclude=false,
            %----------------------------------------------
            footsepline = false,    % Linie über Fußzeilen
            %----------------------------------------------
            headings = normal, % überschriftengröße (Kapitel,...), "big"
            %----------------------------------------------
            fleqn,          % Gleichungen linksbündig statt mittig
            %----------------------------------------------
            %chapterprefix,  % Prefix für Kapitel: "Kapitel"
            %nochapterprefix,% Prefix für Kapitel: keiner.
            %----------------------------------------------
            %halfparskip*,    % Zeilenabstand zwischen Absätzen: mind. 0.25 Zeile
            %halfparskip,    %  Zeilenabstand zwischen Absätzen: 0.5 Zeile
            %parskip,        % Zeilenabstand zwischen Absätzen: 1 Zeile
            %parindent,      % Zeilenabstand zwischen Absätzen: keiner, nur Wort einrücken.
            parskip=half,
            %----------------------------------------------
            german,         % Sprache
            %----------------------------------------------
            %final,          % Version: Endversion
            %draft,          % Version: Zwischenversion, z.~B. ohne Bilder,...
            draft=false,
            %----------------------------------------------
            bibliography = totoc,       % Literaturverzeichnis ins Inhaltsverzeichnis
            %----------------------------------------------
            %pdftex,         % Verwendung mit pdftex
            %----------------------------------------------
            %'BCOR12mm,version=first,       % ergibt links einen Rand von  zusätzlich 1cm zum Binden.
            %----------------------------------------------
            numbers = noenddot,    % Kapitelnummern immer ohne Punkt
            %----------------------------------------------
            caption = tableh,   % korrekte Abstaende bei TabellenUEBERschriften
            %----------------------------------------------
            %abstracton,    % Zusammenfassung?
            %----------------------------------------------
            %draft               % Keine Bilder in der Anzeige, overfull hboxes werden angezeigt
            %----------------------------------------------
                ]{scrbook}[2001/07/30]   % scrbook-Version mind. 2.8j von 2001/07/30
                %]{scrreprt}


% Schalter definieren 
%---------------------------------------------------------
\usepackage{ifthen}
\newboolean{VersionDozent}
\setboolean{VersionDozent}{false}

\newboolean{alphaDIN}
\setboolean{alphaDIN}{true}

% Seitenlayout ändern
% ====================
%\typearea{18}
%\typearea[18mm]{19}          % [Binderand]{Divisor für Seitenaufteilung}
%\areaset[1cm]%               % Zusaetzlicher Rand fuer die Bindung
%        {18cm}{25cm}         % Textbreite und Hoehe


\usepackage{lscape}           % Seite im Querformat bei Erhalt der Kopfzeile
%\usepackage{portland}
\usepackage{rotating} 				% um Text u.s.w. zu rotieren


% Grafiken:
% ==========
\usepackage{ae}                  % Für PDF-Erstellung
\usepackage{graphicx}            % zum Einbinden von Grafiken
\usepackage{float}               % u.a. genaue Plazierung von Gleitobjekten mit H
\graphicspath{{bilder/}{Bilder/}} 			 	 % Verzeichnis, in dem die ganzen Bilder liegen
\usepackage{flafter} 						 % Gleitobjekte immer hinter ihrem ersten Verweis platzieren
%\usepackage{wrapfig} % Bilder mit Text drumrum, Einbinden mit \begin{wrapfigure}
% \usepackage{pstricks}          % PostScript Macros

% pdf-Bilder mit Verwendung von pdf-Latex
%\usepackage[pdftex]{graphicx}
%\usepackage[pdftex,hypertexnames=false]{hyperref}

% Farben vordefinieren
\usepackage{color}
    \definecolor{darkred}{rgb}{0.5,0,0}
    \definecolor{darkgreen}{rgb}{0,0.5,0}
    \definecolor{darkblue}{rgb}{0,0,0.5}
    \definecolor{Hellgrau}{gray}{0.92} 


% Tabellen
% =========
\usepackage{multirow} % Tabellen-Zellen über mehrere Zeilen
\usepackage{ragged2e} %Ausrichtung rechtsbündig möglich mit Option >{\RaggedRight}p{6cm}
\usepackage{longtable} %für 2-Seitige Tabellen

% Für Farben in Tabellen:
\usepackage{colortbl} 



% Bildunterschriften
% ===================
\usepackage{caption, here} % offenbar für neue MikTex-Vers. (z. B. 2.7)


\renewcommand{\captionfont}{\small\itshape} % Bildunterschriften kleiner und kursiv

% Für mehrere Abbildungen in einer:
\usepackage{subfig} %offenbar für neue MikTex-Vers. (z. B. 2.7), funktioniert aber nicht mit den captions (warum???)


% Bei Miktex 2.7 werden die Captions von Subfigures zu breit!
% daher oben subfig laden und hiermit einstellen:
\captionsetup[subfloat]{%
   margin = 4mm,
   labelfont = {it},
   format = plain, % oder 'hang'
   indention = 0em,  % Einruecken der Beschriftung
   singlelinecheck = true, % false (true=bei einer Zeile immer zentrieren)
 }
%}

\captionsetup{
   margin = 4mm,
   %font = {small,rm},
   labelfont = {it},
   format = plain, % oder 'hang'
   indention = 0em,  % Einruecken der Beschriftung
}


% Rechtschreibung 
% ================
\usepackage{ngerman} % neue Rechtschreibung
\usepackage[utf8]{inputenc}    % Input-Encodung: utf-8 fuer Windows
\usepackage{verbatim}          % Quellcode einbinden (\verbatiminput)
\usepackage[T1]{fontenc}         % T1-kodierte Schriften, korrekte Trennmuster fuer Worte mit Umlauten XXX ACHTUNG XXX bringt aber auch mit sich, dass die Auflösung im pdf nicht unendlich ist!!


% Besondere Zeichen
% ==================
% Matrizenbefehle benutzen, für alle Symbole (Mathe)
\usepackage[centertags]{amsmath} % AMS-Mathematik, centertags zentriert Nummer bei split
\usepackage{amsfonts}
\usepackage{amssymb}
\usepackage{amsbsy}

% verschiedene Symbole
%\usepackage{latexsym}           % verschiedene Symbole
\usepackage{textcomp}           % verschiedene Symbole

%für Laplace-Transformationen
\usepackage{trsym, trfsigns}

% für Einheiten
\usepackage{units}

% eingekreiste Zahlen \ding{172} = 1
%\usepackage{psnfss} %wurde nicht gefunden
\usepackage{pifont}

% Häkchen (\checked), Smileys (\smiley),...
%\usepackage{wasysym}

% Kreuz \XSolid, \XSolidBrush
%\usepackage{bbding}

% für z.~B. "N" als Symbol für natürliche Zahlen.
\usepackage{dsfont} %\mathds{N} --> N dann im pdf nicht sichtbar?


% Verzeichnisse, Referenzen
% ==========================
\usepackage{nomencl} 					 % Erstellung einer Nomenklatur
%\usepackage{cite}              % Sortierte und zusammengefasste Zitatnummern 
\usepackage{varioref}          % Verbesserte Referenzen
% \usepackage[numbers, sort]{natbib} % DIN Literaturverzeichnis; nicht zusammen mit cite verwenden!

% Inhaltsverzeichnis: Anzeigetiefe
\setcounter{tocdepth}{3}						%Tiefe des Inhaltsverzeichnisses

% Kopf- und Fußzeilenlayout
% ==========================
\usepackage[left=2cm, right=1.5cm, top=1.5cm, bottom=2cm]{geometry}

\usepackage[automark]{scrlayer-scrpage}
%[] für die erste Seite,{} für alle anderen
\pagestyle{scrheadings}
\clearscrheadfoot                   % Standarddefinitionen l\"{o}schen
\ihead[]{\headmark}                 % Kolumnentitel oben innen
\ohead[\pagemark]{\pagemark}        % Seitennummern immer oben aussen


% BibTex nach DIN1505
\usepackage{bibgerm}    %BibTex German Style DIN 1505

% Seitenlayout
% =============

% damit die erste Zeile eines Absatzes nicht mehr eingerückn wird:
%\parindent0cm

% Werte ändern
\setlength{\parskip}{0.15cm} % Abstand von Absatz zu Absatz (mit leerer Zeile induziert)
\setcounter{secnumdepth}{3} %damit auch Subsubsections nummeriert werden (Nummerierungstiefe)

% Zeilenabstand einstellen:
\usepackage{setspace}            % Zeilenabstand einstellbar
\onehalfspacing                  % eineinhalbzeilig einstellen
%\renewcommand{\baselinestretch}{1.1}

% Das soll weniger Warnungen over/underfull hbox mit sich bringen.
%\sloppy
%\fuzzy  % schaltet \sloppy wieder aus.

% Vordefinierte Namen ändern (Kapitel, ...)
\renewcommand{\bibname}{Literatur} % Literaturverzeichnis wird zu Literatur
\renewcommand{\figurename}{Bild}   % Abbildung wird zu Bild
\renewcommand{\listfigurename}{Bildverzeichnis}


% Schriftart ändern
% ==================
% Serifenlose Schrift als Standard (auch fuer Bildunterschriften!)
\renewcommand{\rmdefault}{cmss}

% Schrift mit Serifen auch fuer Ueberschriften benutzen
%\renewcommand*{\sectfont}{\bfseries}
%\renewcommand*{\descfont}{\bfseries}

\renewcommand{\headfont}{\normalfont\sffamily}    % Kolumnentitel serifenlos
\renewcommand{\pnumfont}{\normalfont\sffamily}    % Seitennummern serifenlos


% Hyperreferenzen -> Links im dvi- und pdf-Dokument
% ==================================================
\usepackage{ifpdf} %ermöglicht das Umschalten von pdf nach eps,...

% Je nach Kompilierung (mit pdf oder mit eps) hier
% automatische Umschaltung.
\ifpdf
    \usepackage[pagebackref,pdftex]{hyperref}
    \pdfcompresslevel=0
\else
    %für Hyperreferenzen im pdf-Dokument, alte Miktex-Version (2.1)
    \usepackage[pagebackref,hypertex]{hyperref}
    % neue Miktex-Version(2.5)
    %\usepackage[a4paper,pagebackref,final,hypertex]{hyperref} %so gehen die Hyperreferenzen in YAP wieder.
\fi

% Hyperreferenzen gezielt einstellen: Farben, Darstellung des pdfs,...
\hypersetup{
    pdftitle        = {Titel},
    pdfsubject      = {Skript},
    pdfauthor       = {Autor},
    pdfkeywords     = {},
    pdfcreator      = {pdftex},
    pdfproducer     = {LaTeX am \today},    linkcolor       = darkred,
    anchorcolor     = darkred,
    citecolor       = darkred,
    filecolor       = darkred,
    menucolor       = darkred,
    urlcolor        = darkred,
%    pdfview         = FitV,
%    pdfstartview    = FitV,
%    pdfstartpage    = {1},
%    pdfpagelayout   =TwoColumnRight,
    pdfpagemode     = UseNone,
    breaklinks      = true,
    colorlinks      = true,
    pdfborder       = 0 0 0,
    plainpages      = false,
%    pdfborder       = 0 0 0,
%    bookmarksnumbered = true,
%    bordercolor     = false,
%    linkbordercolor = {0 0 0},  %Farbe der Textbox
%    citebordercolor = {0 0 0},
%    filebordercolor = {0 0 0},
%    pagebordercolor = {0 0 0},
%    urlbordercolor  = {0 0 0},
%    breaklinks      = false,
%    linktocpage     = true}
%    ps2pdf
}


% Abschluss der Einstellungen
% ============================
\typearea[current]{current}        % Neuberechnung des Satzspiegels mit alten Werten nach Änderung von Zeilenabstand,etc

% Quelltext einbinden
% ============================
\usepackage{listings}
\usepackage{enumitem}
\usepackage{mdwlist}
    %Als Code-Schnippsel im Fließtext
    %\lstinline|print "hello world"|
    %\lstinline{}
    %Als eigenständigen Sourcecode
    %\begin{lstlisting}[caption=Beispielcode]{Name}
    %print "hello world";
    %\end{lstlisting}
    
    %Einbetten einer externen Source-Code-Datei
    %\lstinputlisting[frame=single,label=beispielcode,caption=Ein Beispiel]{beispiel.pl}

%\lstset{numbers=left, numberstyle=\tiny, numbersep=5pt}
%\lstset{language =	Matlab}

\lstset{
  basicstyle=\footnotesize,
	language=C,                % {} für normalen Klartext
	%backgroundcolor=\color{white},
	backgroundcolor=\color{Hellgrau},
	numbers=left, numberstyle=\tiny, numbersep=5pt, %Zeilennummerierung
	linewidth=\linewidth,       % Zeilenbreite
	breaklines=true,             % Zeileumbruch
	breakatwhitespace=true, %Umbruch an Leerzeichen
	%caption = test listing       % Überschrift 
}

% Silbentrennung für schwierige Worte, nur ohne Umlaute
% möglich! Sonst so trennen:  hof-fent-lich
\hyphenation{}

% Nicht deklarierfähige Wörter hier (die sind im Text so zu trennen):
% höhe\-rer